%%%%%%%%%%%%%%%%%%%%%%%%%%%%%%%%%%%%%%%%%%%%%%%%%%%%%%%%%%%%%%%%%%%%%
%
% CSCI 1430 Written Question Template
%
% This is a LaTeX document. LaTeX is a markup language for producing documents.
% You will fill out this document, compile it into a PDF document, then upload the PDF to Gradescope.
%
% To compile it into a PDF you can:
% - Use vscode! Install a LaTeX distribution (all common OS): http://www.latex-project.org/get/
% + VSCode extension: https://marketplace.visualstudio.com/items?itemName=James-Yu.latex-workshop
% - Use an online tool: https://www.overleaf.com/ - most LaTeX packages are pre-installed here (e.g., \usepackage{}).
%
% If you need help with LaTeX, please come to office hours.
% Or, there is plenty of help online:
% https://en.wikibooks.org/wiki/LaTeX
%
% Good luck!
% The CSCI1430 staff
%
%%%%%%%%%%%%%%%%%%%%%%%%%%%%%%%%%%%%%%%%%%%%%%%%%%%%%%%%%%%%%%%%%%%%%

\documentclass{csci1430}

\begin{document}
\title{Homework 5 Written Questions}
\HomeworkShortName{HW5}
\maketitle

\writeinstructions

%%%%%%%%%%%%%%%%%%%%%%%%%%%%%%%%%%%%%%%%%
\pagebreak
\begin{question}[points=14,drawbox=false]
Vision Transformers (ViTs) process images by splitting them into patches, projecting them to latent vectors as patch embeddings, and applying self-attention. In this question, we explore the geometric properties of self-attention and its computational cost.

Recall from HW4 that an operation is \emph{equivariant} to a transformation if transforming the input produces a correspondingly transformed output, and \emph{invariant} if the output does not change at all.

Recall also that we defined a \emph{global} property as one that holds across the whole image equally, and a \emph{local} property as one that holds only within a region.
\end{question}

\begin{subquestion}[points=4]
\textbf{Self-attention without positional encoding---pixel transforms.}
Consider a ViT with no positional encoding. The transformation is applied to the image pixels before patch extraction. Characterize the properties of ``bare'' self-attention. Place \filled~in all relevant answers.
\end{subquestion}

\begin{answer}\noindent
\setlength{\tabcolsep}{12pt}
\begin{tabular}{l|c|c|c|c|c}
\textbf{Transformation} & \rotatebox{90}{\textbf{Global Equivariant}} & \rotatebox{90}{\textbf{Global Invariant}} & \rotatebox{90}{\textbf{Local Equivariant}} & \rotatebox{90}{\textbf{Local Invariant}} & \rotatebox{90}{\textbf{None}} \\
\hline
Translation by $\Delta x=1$ & \cbox & \cbox & \cbox & \cbox & \cbox \\
Translation by $\Delta x=2$ & \cbox & \cbox & \cbox & \cbox & \cbox \\
Rotation by $30^{\circ}$ & \cbox & \cbox & \cbox & \cbox & \cbox \\
Rotation by $90^{\circ}$ & \cbox & \cbox & \cbox & \cbox & \cbox \\
Scale by $2.0\times$ & \cbox & \cbox & \cbox & \cbox & \cbox \\
Scale by $0.3\times$ & \cbox & \cbox & \cbox & \cbox & \cbox \\
Reflection (horizontal flip) & \cbox & \cbox & \cbox & \cbox & \cbox \\
Permutation (pixel pos.~swap) & \cbox & \cbox & \cbox & \cbox & \cbox \\
\end{tabular}
\end{answer}

\pagebreak
\begin{subquestion}[points=2]
\textbf{Self-attention without positional encoding---patch embeddings.}
Beyond the pixel level to consider the patch embeddings themselves, what invariances or equivariances hold through a self attention layer? Why? \textbf{[1--2 sentences]}
\end{subquestion}

\begin{answer}[height=14]

\end{answer}

\pagebreak
\begin{subquestion}[points=4]
\textbf{Adding positional encoding.}

ViTs add a learned positional encoding to each patch embedding before self-attention. Each spatial position in the patch grid has a unique, fixed encoding vector added to its embedding.

How does this change for each of translation, rotation, scaling, reflection, and permutation properties from Q1.1 and Q1.2? \textbf{[4--8 sentences]}
\end{subquestion}

\begin{answer}[height=28]

\end{answer}


\pagebreak
\begin{subquestion}[points=4]
\textbf{Computational cost of self-attention.}

Self-attention computes $\text{softmax}(QK^\top / \sqrt{d})V$, where $Q, K, V \in \mathbb{R}^{N \times D}$, $N$ is the number of patches, and $D$ is the embedding dimension. The dominant cost is the $QK^\top$ matrix multiplication, which is $O(N^2 D)$.

For a $224 \times 224$ image with $16 \times 16$ patches, $N = 14 \times 14 = 196$.

\vspace{0.5em}
\noindent (a) If we halve the patch size to $8 \times 8$, how many patches $N$ do we get?

\vspace{0.5em}
\noindent (b) By what factor does the self-attention cost change compared to $16 \times 16$ patches?

\vspace{0.5em}
\noindent (c) By contrast, if the number of spatial positions in a feature map quadruples (as it does when halving patch size), by what factor does the cost of a standard convolutional layer change?

\vspace{0.5em}
\noindent (d) Using larger patch sizes ameliorates the computational cost of self-attention. What trade-off does varying patch size induce when considering the visual features within images? \textbf{[2--3 sentences]}
\end{subquestion}

\begin{answer}[height=20]
(a) $N =$

\vspace{1em}
(b) Factor:

\vspace{1em}
(c) CNN factor:

\vspace{1em}
(d)
\end{answer}


%%%%%%%%%%%%%%%%%%%%%%%%%%%%%%%%%%%%%%%%%%%%%
\pagebreak

\begin{question}[points=12,drawbox=false]
DINO is a self-supervised optimization method that learns a representation by having a student network match the output of a momentum teacher network across different views of the same image. In this question, we explore how DINO's hyperparameters affect training and the learned representation.
\end{question}

\begin{subquestion}[points=3]
\textbf{Temperature-sharpened softmax.}

In DINO, both the student and teacher produce output vectors $\mathbf{z} \in \mathbb{R}^K$. These are converted to probability distributions using a \emph{temperature-sharpened softmax}:
\[
p_i = \frac{\exp(z_i / \tau)}{\sum_{j=1}^{K} \exp(z_j / \tau)}
\]
where $\tau > 0$ is the temperature. 

What is the effect of \emph{decreasing} $\tau$ toward zero? Select all that apply.
\end{subquestion}

\begin{answerlist}
    \item The output distribution becomes more uniform (all $p_i$ become equal).
    \item The output distribution becomes more peaked (one $p_i$ dominates).
    \item The logits are scaled but their relative ordering is unchanged, so the relative output distribution stays the same.
    \item The entropy of the output distribution increases.
    \item The entropy of the output distribution stays the same.
    \item The entropy of the output distribution decreases.

\end{answerlist}


\pagebreak
\begin{subquestion}[points=3]
\textbf{Temperature asymmetry.}

DINO uses a lower temperature for the teacher ($\tau_t = 0.04$) than the student ($\tau_s = 0.1$).

What happens to the teacher's output distribution as $\tau_t \rightarrow 0$? What about as $\tau_t \rightarrow \infty$? Why is it important that $\tau_t < \tau_s$? \textbf{[3--4 sentences]}
\end{subquestion}

\begin{answer}[height=14]

\end{answer}


\pagebreak
\begin{subquestion}[points=3]
\textbf{EMA momentum and training stability.}

The teacher's parameters $\theta_t$ are not trained by gradient descent. Instead, they are updated each step via an \emph{exponential moving average} (EMA) of the student's parameters $\theta_s$:
\[
\theta_t \leftarrow \lambda \cdot \theta_t + (1 - \lambda) \cdot \theta_s
\]
where $\lambda \in [0, 1]$ is the \emph{momentum}. DINO uses $\lambda = 0.996$.

Describe how the teacher's targets evolve and whether training succeeds or fails in each of these three scenarios. \textbf{[2 sentences each]}
\end{subquestion}

\begin{answer}[height=16]
(a) $\lambda = 0$ (teacher is set equal to student every step):

\vspace{2.5em}
(b) $\lambda = 1$ (teacher never updates):

\vspace{2.5em}
(c) $\lambda = 0.996$ (DINO's default):
\end{answer}


\pagebreak
\begin{subquestion}[points=3]
\textbf{Cross-view pairing.}

The teacher sees two global crops---large views covering 40--100\% of the image. The student sees all crops---two global + several local crops covering 5--40\% of the image. All crops are placed into a shared array. The loss is computed over all \emph{cross-view} pairs: teacher crop $i$ with student crop $j$, where $i \neq j$.

Why do we skip same-view pairs ($i = j$)? And why does the teacher see only global crops while the student also sees local crops? \textbf{[4--6 sentences]}
\end{subquestion}

\begin{answer}[height=20]

\end{answer}


%%%%%%%%%%%%%%%%%%%%%%%%%%%%%%%%%%%
\pagebreak

\writefeedback

% \pagebreak
% \section*{Any additional pages would go here.}


\end{document}
